\begin{frame}{자주 묻는 질문(15.2, 이어서)}
  \small
  \begin{itemize}
    \item \textbf{Q: ``승률 98.7\%''는 무작위 상대와의 실전
      승률(67\%)과 같은 의미인가?} --- \textbf{아니} --- 이
      구분이 ``트리와 롤아웃의 역할 분담'' 절의 핵심.
      98.7\% = \textbf{트리가 ``상대 대응에 대한 내 필승
      응답''을 통계에 저장한 뒤}의 \textbf{합성 추정치},
      순수 무작위 승률($2/3$)을 \textbf{뛰어넘음}.
      시뮬레이션을 늘리면 님처럼 \textbf{끝이 짧은} 게임에서는
      100\%(최선-최선 기준 참 승률)에, 바둑처럼 \textbf{긴}
      게임에서는 ``무작위 참승률''과 ``최선 상한''
      \textbf{사이에} 수렴 --- 그 위치가 바로 \textbf{신경망
      롤아웃이 개입}하는 곳
    \item \textbf{Q: 시뮬레이션은 몇 번이면 충분한가?} ---
      \textbf{닫힌 형태의 공식이 없음} ---
      (1) \textbf{수렴 곡선이 평탄해지나}(이기는 수의 방문
      점유율이 고정값에 다가서는가), (2) \textbf{수의 순위}가
      시뮬레이션 추가에 대해 \textbf{변하지 않는가}로
      실측해 판단. 예산 배분 경험칙 = ``\textbf{중요한 수
      (결정적 국면)에, 쓸 수 있는 만큼}'' --- AlphaGo는 한
      수당 CPU 수 초$\sim$수 분, 남은 깊이가 짧은 \textbf{끝물}은
      수렴이 훨씬 빨라 같은 횟수로 더 정확한 추정
  \end{itemize}
\end{frame}
